\begin{table}[tb]
\centering\small\setlength{\tabcolsep}{4pt}
\caption{同条件の反復間12対．一致数／使用量既知の判定可能対数と，判定可能数／全４対を示す．}
\begin{tabular}{lrrr}
条件 & 判定可能 & 全要求一致 & 全公開出力一致 \\ \hline
M & 4/4 & 0/4 & 0/4 \\
H & 3/4 & 3/3 & 3/3 \\
L & 3/4 & 3/3 & 3/3 \\
\hline\end{tabular}
\end{table}
\begin{table}[tb]
\centering\small\setlength{\tabcolsep}{4pt}
\caption{全８件のLで最初に実測残時間を通知した生成番号．未到達は効果がないことの証明ではない．}
\begin{tabular}{lrr}
課題 & 反復1 & 反復2 \\ \hline
写真移動 & 未到達 & 未到達 \\
精算 & 未到達 & 未到達 \\
返金 & 未到達 & 未到達 \\
曲情報の書出し & 16 & 16 \\
\hline\end{tabular}
\end{table}
反復間12対に加え，反復内の条件間24対を全件保存した．これら36対は同じ実行を共有し，独立標本ではない．現在の通知だけを揃えた補助照合は，過去の履歴差を消去せず，モデルtoken列や全仕事費用の因果効果を同定しない．
